\section{Ley de Newton del enfriamiento}

La ley de Newton describe el intercambio de calor de un cuerpo con su entorno cuando domina la convección y la diferencia de temperaturas es moderada: la tasa de variación de la temperatura es proporcional a dicha diferencia.

\begin{modelo}[enfriamiento de Newton]
Sean $T(t)$ la temperatura del cuerpo y $S(t)$ la del entorno. Entonces
\begin{equation}
  \dv{T}{t} = -k\,\bigl(T - S\bigr), \qquad T(0)=T_0,
  \label{eq:newton}
\end{equation}
con $k>0$ constante de enfriamiento, $[k]=\si{\per\second}$.
\end{modelo}

\begin{remark}
El signo de \eqref{eq:newton} expresa que el calor fluye del cuerpo más caliente al más frío: si $T>S$ entonces $\dot T<0$, y si $T<S$ entonces $\dot T>0$.
\end{remark}

\subsection*{Entorno a temperatura constante}

Si el entorno es un foco térmico ideal, $S(t)\equiv S_0$, y \eqref{eq:newton} es lineal de coeficientes constantes: $\dot T + kT = kS_0$.

\begin{proposition}[solución del problema de Newton]\label{prop:newton}
La solución del problema \eqref{eq:newton} con $S\equiv S_0$ es
\begin{equation}
  T(t) = S_0 + (T_0 - S_0)\,\ee^{-kt},
  \label{eq:newton-sol}
\end{equation}
y $T(t)\to S_0$ cuando $t\to\infty$.
\end{proposition}

\begin{proof}
La ecuación homogénea $\dot T + kT=0$ tiene solución general $T_h = C\,\ee^{-kt}$, y $T_p\equiv S_0$ es una solución particular. Luego $T = S_0 + C\,\ee^{-kt}$, y la condición inicial $T(0)=T_0$ da $C=T_0-S_0$. Como $k>0$, el término exponencial tiende a cero.
\end{proof}

\begin{remark}
La magnitud $\tau = 1/k$ es la \emph{constante de tiempo}: la diferencia $T-S_0$ se reduce en un factor $\ee^{-1}$ cada $\tau$.
\end{remark}

\begin{figure}[H]
  \centering
  \begin{tikzpicture}
    \begin{axis}[informe, xmin=0, xmax=120, ymin=0, ymax=120,
      xlabel={$t$ (\si{\minute})}, ylabel={$T$ (\si{\celsius})},
      legend style={at={(0.97,0.62)}}]
      \addplot[asintota, forget plot, domain=0:120, samples=2] {20};
      \node[font=\small, text=gris, anchor=south east] at (axis cs:118,21) {$S_0$};
      \addplot[serie1, domain=0:120, samples=100] {20 + 90*exp(-0.04*x)};
      \addlegendentry{$T_0>S_0$}
      \addplot[serie2, domain=0:120, samples=100] {20 - 15*exp(-0.04*x)};
      \addlegendentry{$T_0<S_0$}
    \end{axis}
  \end{tikzpicture}
  \caption{Ley de Newton con $S_0=\SI{20}{\celsius}$ y $k=\SI{0,04}{\per\minute}$. Toda solución converge a $S_0$.}
  \label{fig:newton}
\end{figure}

\begin{example}
Un cuerpo a $T_0=\SI{90}{\celsius}$ se sitúa en un entorno a $S_0=\SI{20}{\celsius}$. Tras \SI{10}{\minute} su temperatura es \SI{60}{\celsius}. Determinar $k$ y el instante en que alcanza \SI{30}{\celsius}.

De \eqref{eq:newton-sol}, $60 = 20 + 70\,\ee^{-10k}$, luego
\[
  k = \tfrac{1}{10}\ln\tfrac{70}{40} \approx \SI{0,056}{\per\minute}.
\]
Imponiendo $T(t)=30$: $\ee^{-kt}=\tfrac{10}{70}$, de donde
$t = \tfrac{1}{k}\ln 7 \approx \SI{34,8}{\minute}$.\finejemplo
\end{example}

\newpage
